\subsection{Accuracy, optimizer memory, and follow-up experiments}
\label{sec:strength-followup}
These experiments measure prediction quality and optimizer-state memory
together, then test subspace tracking on controlled data streams. They also
repeat the PINN and graph comparisons with larger validation budgets.
All comparisons, including those favoring another method, are reported below.

\paragraph{Prediction quality and optimizer memory.}
On PubMed the original extension obtains 76.0\% with GrassRoute versus
74.2\% with AdamW, retaining 87,576 versus 1,031,216 state bytes: a 91.5\%
reduction and a 1.8 percentage-point accuracy gain in the single paired run.
The periodic implementation is more accurate at 77.1\% with the same state
cost. On all-user MovieLens, GrassRoute retains 25.3\% of AdamW's state and
obtains 12.247\% versus 12.363\% Recall@20, a 0.116 percentage-point loss.
This memory--quality trade-off is descriptive; significance and
noninferiority were not tested. On Poisson its relative solution error is
$1.255\times10^{-4}$ versus $2.987\times10^{-4}$ for the original periodic
implementation at the same persistent state cost, but AdamW is more accurate.
Table~\ref{tab:state-quality} reports every completed paired task,
including regression and operator-learning results that favor the baselines.

\paragraph{Tracking with the implemented optimizers.}
We run the routed optimizer code on Gaussian gradient matrices with a known
rotating population covariance, using rank 4 in dimension 64, signal/tail
eigenvalues $1$ and $0.03$, and 16 or 128 gradient rows per update. The four
rotation rates are $0,0.001,0.005,0.02$ radians per step. All methods start
from the same random frame and observe the same 2000 matrices. We compare
GrassRoute at a fixed $\rho=3$, its no-momentum ablation, a fixed router,
and transport-corrected instantaneous SVD refreshes at periods 5, 25 and 100.
The target projector is known analytically. We measure error
\emph{before} each update so that an SVD refresh is not evaluated only
on the batch used to compute it.
With the fixed tracking rate $\rho=3$, GrassRoute has lower steady tracking error than interval-25 instantaneous SVD in 8 of the 8 predeclared conditions. Table~\ref{tab:tracking-strength} includes the static and fast-drift controls, the no-momentum ablation, and interval-5 and interval-100 alternatives. These comparisons measure subspace identification on synthetic streams, not an end-to-end accuracy advantage.

% WWW table style: three rules, bold headers, reference green highlight.
\begin{table*}[t]\centering\footnotesize
\caption{\textbf{Tracking implementation resources.} Mean CPU optimizer-call time across all four drift rates at 128 gradient rows, including diagnostics; persistent state is decimal kB. Shared implementation timing is reported separately from task accuracy.}
\label{tab:tracking-resources}
\setlength{\tabcolsep}{5pt}
\reviewfit{\begin{tabular}{lrr}
\toprule {\bfseries Method} & {\bfseries CPU call / ms} & {\bfseries State / kB}\\\midrule
\rowcolor{oursbg}\textbf{GrassRoute} &  0.1497  &  6.14\\
NoMomentum  &  0.1354  &  6.14\\
Fixed  &  0.1428  &  6.14\\
Periodic5  &  0.09495  &  6.14\\
Periodic25  &  0.0632  &  6.14\\
Periodic100  &  0.05768  &  6.14\\
\bottomrule\end{tabular}}
\end{table*}

\IfFileExists{figures/strength_tracking_table.tex}{% WWW table style: three rules, bold headers, reference green highlight.
\begin{table*}[htbp]\centering\scriptsize
\caption{Complete predeclared tracking grid, 2000 steps, seed 42. Mean $r^{-1}\|\sin\Theta\|_F^2$ before each update on steps 1501--2000; lower is better. All methods see identical streams and start from the same random frame. Periodic variants use corrected diagonal transport. Bold marks the minimum among all displayed methods.}
\label{tab:tracking-strength}
\reviewfit{\begin{tabular}{rrrrrrrr}\toprule {\bfseries Rows} & {\bfseries Drift} & \cellcolor{oursbg}{\bfseries GrassRoute} & {\bfseries No mom.} & {\bfseries Fixed} & {\bfseries Periodic5} & {\bfseries Periodic25} & {\bfseries Periodic100}\\\midrule
16  &  0.000  &  \cellcolor{oursbg}\textbf{0.0307} &  0.0410  &  0.9119  &  0.1455  &  0.1386  &  0.1403\\
16  &  0.001  &  \cellcolor{oursbg}\textbf{0.0307} &  0.0410  &  0.9346  &  0.1455  &  0.1388  &  0.1429\\
16  &  0.005  &  \cellcolor{oursbg}\textbf{0.0308} &  0.0411  &  0.9223  &  0.1456  &  0.1435  &  0.2070\\
16  &  0.020  &  \cellcolor{oursbg}\textbf{0.0322} &  0.0427  &  0.9263  &  0.1489  &  0.2109  &  0.6477\\
128  &  0.000  &  \cellcolor{oursbg}\textbf{0.0038} &  0.0050  &  0.9119  &  0.0156  &  0.0161  &  0.0169\\
128  &  0.001  &  \cellcolor{oursbg}\textbf{0.0038} &  0.0050  &  0.9346  &  0.0157  &  0.0164  &  0.0205\\
128  &  0.005  &  \cellcolor{oursbg}\textbf{0.0039} &  0.0051  &  0.9223  &  0.0159  &  0.0216  &  0.0973\\
128  &  0.020  &  \cellcolor{oursbg}\textbf{0.0054} &  0.0067  &  0.9263  &  0.0201  &  0.0988  &  0.6060\\
\bottomrule\end{tabular}}\end{table*}
}{}

\paragraph{Covariance smoothing and faster drift.}
An instantaneous SVD does not average sample noise across steps. To test
whether averaging alone explains the gain, a separate protocol adds
normalized full-covariance EMA estimators with decays $0.9$ and $0.99$,
refreshed every 5 or 25 steps. It repeats every original condition and adds
rotation rate $0.1$ at both row counts. The estimators store the full
$64\times64$ covariance, which is included in the state cost. They use
the same moment transport, initial frame, gradient stream, and pre-update
metric. This control provides temporal smoothing without Grassmann momentum,
at the cost of additional memory.
\IfFileExists{figures/strength_tracking_ema_summary.tex}{After adding four full-covariance EMA controls and the nonadiabatic rate $0.1$, GrassRoute has the lowest error in 5 of 10 conditions and beats both instantaneous refresh baselines in 10 of 10. EMA SVD is stronger in the static or slow-drift conditions. At 128 rows and drift $0.02$, GrassRoute has mean error 0.00541 versus 0.05440 for the best displayed EMA alternative, with 6,144 versus 22,528 persistent state bytes. This isolates a noise--lag--memory trade-off; the current implementation does not establish a CPU speed advantage.
}{}

\IfFileExists{figures/strength_tracking_ema_table.tex}{% WWW table style: three rules, bold headers, reference green highlight.
\begin{table*}[htbp]\centering\scriptsize\setlength{\tabcolsep}{3pt}
\caption{Stronger tracking controls on all ten predeclared conditions. EMA SVD stores the full normalized covariance and uses decay .9 or .99; its refresh period is 5 or 25. Mean pre-update $r^{-1}\|\sin\Theta\|_F^2$ on steps 1501--2000. Bold identifies the best displayed error; no temporal observations are treated as independent replicates.}
\label{tab:tracking-ema}
\reviewfit{\begin{tabular}{rrrrrrrrr}\toprule {\bfseries Rows} & {\bfseries Drift} & \cellcolor{oursbg}{\bfseries GrassRoute} & {\bfseries SVD/5} & {\bfseries SVD/25} & {\bfseries EMA.9/5} & {\bfseries EMA.9/25} & {\bfseries EMA.99/5} & {\bfseries EMA.99/25}\\\midrule
16  &  0.000  &  \cellcolor{oursbg}0.03073 &  0.14550  &  0.13859  &  0.00636  &  0.00644  &  \textbf{0.00059}  &  0.00059\\
16  &  0.001  &  \cellcolor{oursbg}0.03074 &  0.14549  &  0.13882  &  \textbf{0.00649}  &  0.00698  &  0.01054  &  0.01269\\
16  &  0.005  &  \cellcolor{oursbg}0.03083 &  0.14565  &  0.14349  &  \textbf{0.00996}  &  0.01974  &  0.15492  &  0.19335\\
16  &  0.020  &  \cellcolor{oursbg}\textbf{0.03217} &  0.14885  &  0.21093  &  0.06017  &  0.19458  &  0.42936  &  0.62078\\
16  &  0.100  &  \cellcolor{oursbg}\textbf{0.07407} &  0.23144  &  0.65302  &  0.51743  &  0.60934  &  0.63732  &  0.54925\\
128  &  0.000  &  \cellcolor{oursbg}0.00381 &  0.01564  &  0.01613  &  0.00078  &  0.00077  &  0.00007  &  \textbf{0.00007}\\
128  &  0.001  &  \cellcolor{oursbg}0.00381 &  0.01566  &  0.01635  &  \textbf{0.00093}  &  0.00134  &  0.01021  &  0.01235\\
128  &  0.005  &  \cellcolor{oursbg}\textbf{0.00391} &  0.01594  &  0.02159  &  0.00445  &  0.01420  &  0.15473  &  0.19328\\
128  &  0.020  &  \cellcolor{oursbg}\textbf{0.00541} &  0.02008  &  0.09878  &  0.05440  &  0.18954  &  0.42786  &  0.62058\\
128  &  0.100  &  \cellcolor{oursbg}\textbf{0.04949} &  0.11814  &  0.60853  &  0.50970  &  0.60812  &  0.70552  &  0.57523\\
\bottomrule\end{tabular}}\end{table*}
}{}

\paragraph{Larger validation search on exact-solution PINNs.}
Motivated by the archived Burgers result, we evaluate both exact-solution
PINNs under a separate protocol. Each of AdamW, fixed routing,
\textsc{PeriodicDiag}, and GrassRoute receives nine 100-epoch validation
screening runs. AdamW searches nine geometrically spaced learning rates from
$10^{-4}$ to $0.03$. Each routed method searches learning rates
$\{0.001,0.003,0.01\}$, crossed with a method-specific parameter:
residual-rate ratios $\{0.3,1,3\}$ for fixed routing; refresh periods
$\{5,25,100\}$ for \textsc{PeriodicDiag}; and tracking rates
$\{0.3,1,3\}$ for GrassRoute. Every candidate and divergence is recorded. Final validation relative
solution error selects the setting; a fresh model then trains for 2000 full
epochs and is tested once. Methods share parameter initialization and
collocation streams. The protocol and source hash were saved before
training. Because the search budget and periodic moment transport differ,
these results are separate from the 162-comparison extension.
The independent follow-up has completed 8 of 8 selected formal runs. The currently completed follow-up endpoints do not establish an all-baseline GrassRoute win. Its endpoints are not pooled with the original six-method extension.

On Poisson, GrassRoute's selected relative solution error is
$6.70\times10^{-5}$ at 0.0379 MB of persistent optimizer state, versus
$9.52\times10^{-4}$ at 0.2673 MB for the independently selected AdamW
configuration. The comparison depends on selection: the original, smaller search found
an AdamW error of $9.28\times10^{-5}$, while \textsc{PeriodicDiag} in
the new search reaches $4.90\times10^{-5}$. AdamW also remains better on
forced Burgers. Screening for 100 epochs can select a setting that performs
poorly over the full horizon, even with more candidates. The table below
reports all selected results.
\IfFileExists{figures/strength_pinn_table.tex}{% WWW table style: three rules, bold headers, reference green highlight.
\begin{table*}[htbp]\centering\small
\caption{Independent expanded-search PINN follow-up with a transport-corrected periodic comparator. Nine validation-screen candidates per method (100 epochs each), then a fresh 2000-epoch run of the selected configuration. Final held-out solution relative $L_2$ error; lower is better. -- denotes unfinished, not omitted, training.}
\label{tab:pinn-strength}
\reviewfit{\begin{tabular}{lrrrr}\toprule {\bfseries Task} & {\bfseries AdamW} & {\bfseries Fixed} & {\bfseries PeriodicDiag} & \cellcolor{oursbg}{\bfseries GrassRoute}\\\midrule
Poisson / PINN  &  $9.52\!\times\!10^{-4}$  &  $1.97\!\times\!10^{-4}$  &  $\mathbf{4.90\!\times\!10^{-5}}$  &  \cellcolor{oursbg}$6.70\!\times\!10^{-5}$\\
Forced Burgers / PINN  &  \textbf{0.00239}  &  0.0128  &  0.00582  &  \cellcolor{oursbg}0.00427\\
\bottomrule\end{tabular}}\end{table*}
}{}

\paragraph{Expanded graph search.}
A separate follow-up tests Cora, CiteSeer and PubMed against all six methods,
including corrected periodic transport and official SubTrack++. Each method
receives nine 100-epoch validation candidates and 1000 fresh formal epochs.
GrassRoute crosses three learning rates with tracking rates
$\{0.3,1,3\}$; fixed routing searches residual-rate ratios, corrected periodic
routing searches refresh intervals, and official SubTrack++ searches its
initial geodesic step. AdamW and SGD each search nine learning rates. The protocol specifies every candidate before training and uses paired
model and dropout streams with the public graph splits. Every selected
result is retained. These transductive experiments test optimizer
performance; they do not verify the generalization theorem's iid assumptions.
Numerical failures, bounded process retries, and the CPU device adapter are
documented in Appendix~\ref{app:numerical-recovery}. They preserve the
nine unique hyperparameters, and fallback choices never use test scores.
The independent Planetoid follow-up has completed 18 of 18 selected formal runs. Completed paired scores favor GrassRoute over both AdamW and corrected periodic routing on Cora / SAGE. Pending cells and full selection/state records are retained.

